\documentclass[conference]{IEEEtran}
\IEEEoverridecommandlockouts

\usepackage{cite}
\usepackage{amsmath,amssymb,amsfonts}
\usepackage{graphicx}
\usepackage{textcomp}
\usepackage{xcolor}
\usepackage{booktabs}
\usepackage{multirow}
\usepackage{microtype}

\def\BibTeX{{\rm B\kern-.05em{\sc i\kern-.025em b}\kern-.08em
    T\kern-.1667em\lower.7ex\hbox{E}\kern-.125emX}}

\begin{document}

\title{MinImageScorer: Error-Driven Difficulty Scoring\\
for Guided Copy-Paste Augmentation in\\
Agricultural Object Detection}

% TODO: Fill in your name, supervisor name, affiliation
\author{
\IEEEauthorblockN{Author Name\textsuperscript{1},
                  Supervisor Name\textsuperscript{1}}
\IEEEauthorblockA{%
\textsuperscript{1}Department of XXX, University of XXX, Country\\
email@university.edu}
}

\maketitle

% ----------------------------------------------------------
\begin{abstract}
Copy-paste augmentation is widely applied in object detection
but is rarely evaluated against the specific failure mode a
dataset exhibits. We propose \textbf{MinImageScorer}, an
error-driven framework that scores per-object detection difficulty
from baseline model failures, combining confidence and spatial
accuracy penalties calibrated via Pearson correlation optimization.
We validate the score on two agricultural benchmarks---MinneApple
and Global Wheat Head (GWHD)---demonstrating that it reliably
ranks objects and images by difficulty (Pearson $r$ up to $0.82$
with miss rate) and correctly recovers domain- and size-level
failure patterns without supervision. We then use the score to
guide copy-paste augmentation in a controlled three-seed study.
Score-guided selection outperforms random selection on MinneApple
($\Delta\mathrm{AP}{=}{+}0.013$) but neither condition recovers
the no-augmentation baseline. Feature-space analysis confirms
that augmented images shift the training distribution away from
the test set, explaining the degradation. These findings establish
that MinImageScorer correctly identifies hard objects, and that
copy-paste utility depends on the dataset's underlying failure mode
rather than on selection strategy alone.
\end{abstract}

\begin{IEEEkeywords}
object detection, copy-paste augmentation, difficulty scoring,
agricultural detection, small objects, domain shift
\end{IEEEkeywords}

% ----------------------------------------------------------
\section{Introduction}

Agricultural object detection benchmarks such as
MinneApple~\cite{hani2020minneapple} and Global Wheat Head
(GWHD)~\cite{david2021global} present distinct structural failure
modes. MinneApple is dominated by high-density small objects near the
resolution floor of standard backbones, while GWHD exhibits
substantial cross-domain visual variation across geographic sources.
Copy-paste augmentation has been shown to improve small-object AP on
MS COCO~\cite{kisantal2019augmentation,ghiasi2021simple}, but those
datasets lack structural resolution limits or domain shift---the
conditions that dominate agricultural benchmarks.

A critical gap remains: practitioners apply copy-paste without first
diagnosing whether the dataset's failure mode is addressable by adding
more instances. When failure is structural, copy-paste may corrupt the
training distribution rather than improve it. No prior work provides
an empirical scoring framework to diagnose this mismatch. We make
three contributions:
\begin{enumerate}
    \item \textbf{MinImageScorer}: a per-object difficulty score
    derived from model failures, with data-driven weight calibration
    requiring no domain assumptions.
    \item \textbf{Empirical score validation}: the score reliably
    correlates with miss rate and false-positive rate, and correctly
    recovers domain-level and object-size-level difficulty rankings
    without supervision.
    \item \textbf{Controlled augmentation study}: score-guided
    copy-paste is compared to random selection across three seeds,
    with feature-space analysis explaining the outcome.
\end{enumerate}

% ----------------------------------------------------------
\section{Methodology}

\subsection{MinImageScorer Formulation}

Let $G(I)$ be the ground-truth objects in image $I$ and
$\mathcal{P}(g;\tau)$ the model predictions satisfying
$\mathrm{IoU}(p,g)\ge\tau$. The \emph{object difficulty score} is:

\begin{equation}
S_{\mathrm{obj}}(g) =
\begin{cases}
\displaystyle\min_{p\,\in\,\mathcal{P}(g;\tau)}
  \bigl[
    \alpha\,(1 - \mathrm{conf}(p))
    + \beta\,(1 - \mathrm{IoU}(p,g))
  \bigr],
  & \mathcal{P}(g;\tau) \neq \varnothing, \\[4pt]
\alpha + \beta,
  & \text{otherwise.}
\end{cases}
\label{eq:obj}
\end{equation}

Missed detections receive the maximum penalty $\alpha{+}\beta$.
The \emph{image-level score} aggregates object difficulty and
penalizes false positives:

\begin{equation}
S_{\mathrm{img}}(I)
  = w_1 \cdot \frac{1}{|G(I)|}\sum_{g\in G(I)} S_{\mathrm{obj}}(g)
  \;+\; w_2 \cdot \mathrm{FPrate}(I).
\label{eq:img}
\end{equation}

We use two prediction sets: \textit{low-confidence} predictions
(dense, $\tau{=}0.01$) for $S_{\mathrm{obj}}$, capturing near-misses
and weak detections; and \textit{optimal-threshold} predictions for
$\mathrm{FPrate}$, reflecting deployment-level errors.
Weights $\alpha,\beta,w_1,w_2$ are calibrated to maximize Pearson
correlation with per-image miss rate and FP rate on the training set.

\subsection{Score-Guided Copy-Paste Pipeline}

(1)~Train a baseline YOLOv8s~\cite{yolov8} with all built-in
augmentations disabled. (2)~Run inference to collect confidence and
IoU signals. (3)~Compute $S_{\mathrm{obj}}$ and $S_{\mathrm{img}}$.
(4)~For score-guided augmentation, sample objects proportionally to
$S_{\mathrm{obj}}$ via exponential weighting; for random, sample
uniformly. (5)~Paste 2--8 objects per augmented image (within-image
only) and retrain. Both conditions are identical except for the
selection policy.

% ----------------------------------------------------------
\section{Score Validation}

Before applying the score to guide augmentation, we validate that it
reliably measures detection difficulty across two structural axes:
per-image failure correlation and per-domain/object-size rankings.

\subsection{Correlation with Model Failures}

Table~\ref{tab:corr} reports Pearson correlation between
$S_{\mathrm{img}}$ and per-image miss rate / FP rate, computed on
training images across both datasets. All cases show strong positive
correlation. The strongest relationship is on GWHD with traditional
augmentation ($r{=}0.82$ vs.\ miss rate, $r{=}0.82$ vs.\ FP rate),
where the baseline model is stable and the score signal is
clean. On MinneApple without traditional augmentation,
stability drops (mean inter-seed Pearson ${\approx}0.67$), because
unstabilized models produce fluctuating confidence scores on
densely clustered small objects. This confirms that score reliability
depends on baseline model stability, not on the score design itself.

% ----------------------------------------------------------------
% FIGURE 1 RECOMMENDATION:
% Replace this placeholder with a 2-panel scatter plot:
%   Left panel:  x = S_img (MinneApple), y = miss rate per image.
%   Right panel: x = S_img (GWHD),        y = miss rate per image.
% Each point is one training image. Fit a regression line.
% This is your strongest single visualization for Claim 1.
% Export as score_correlation.pdf (width = \columnwidth).
% ----------------------------------------------------------------
\begin{figure}[t]
\centering
%\includegraphics[width=\columnwidth]{score_correlation.pdf}
\framebox{\parbox{0.88\columnwidth}{\centering\vspace{1.5cm}
\small FIGURE 1 PLACEHOLDER\\
Scatter: $S_{\mathrm{img}}$ vs.\ per-image miss rate\\
Left: MinneApple \quad Right: GWHD\\
Each point = one training image; regression line shown.
\vspace{1.5cm}}}
\caption{Score vs.\ per-image miss rate on training images.
  MinneApple (Pearson $r{=}0.70$); GWHD ($r{=}0.82$, w/ trad.\ aug).
  Higher score correctly predicts harder images on both datasets.}
\label{fig:corr}
\end{figure}

\begin{table}[t]
\centering
\caption{Pearson correlation of $S_{\mathrm{img}}$ with per-image
miss rate and FP rate (training set, w/ traditional augmentation).}
\label{tab:corr}
\resizebox{\columnwidth}{!}{%
\begin{tabular}{@{}lcccc@{}}
\toprule
\textbf{Dataset}
  & $r$(\,score,\,miss\,rate\,)
  & $r$(\,score,\,FP\,rate\,)
  & Inter-seed $r$ (mean) \\
\midrule
MinneApple  & 0.698 & 0.688 & 0.806 \\
GWHD        & 0.816 & 0.818 & 0.900 \\
\bottomrule
\end{tabular}%
}
\end{table}

\subsection{Domain- and Size-Level Validity}

Beyond per-image correlation, we verify that the score recovers
structural difficulty known from ground-truth AP. On GWHD,
Table~\ref{tab:domain} shows that domains with lower AP receive
higher mean scores: the \textit{arvalis~2} domain (AP\,$=\,0.365$)
receives a mean score of $0.431$, while the easiest domain
\textit{rres~1} (AP\,$=\,0.572$) scores only $0.260$.
This domain ranking is recovered \textit{without access to test
labels}---purely from training-set model failures.

On MinneApple, small objects (${\le}32$px) average a score of
$0.378$ vs.\ $0.222$ for medium objects, perfectly matching the
AP gap ($0.132$ vs.\ $0.522$). Together, these results confirm
that MinImageScorer captures both image-level and structural
difficulty without domain-specific tuning.

\begin{table}[t]
\centering
\caption{GWHD domain AP vs.\ mean $S_{\mathrm{img}}$. Domains with
  lower AP consistently receive higher scores (no test labels used).}
\label{tab:domain}
\resizebox{\columnwidth}{!}{%
\begin{tabular}{@{}lrrrr@{}}
\toprule
\textbf{Domain} & \textbf{n\,img} & \textbf{AP} & \textbf{AP\textsubscript{50}} & \textbf{Mean Score} \\
\midrule
rres~1      & 363 & 0.572 & 0.962 & 0.260 \\
inrae~1     & 146 & 0.523 & 0.954 & 0.284 \\
ethz~1      & 592 & 0.550 & 0.940 & 0.291 \\
arvalis~3   & 457 & 0.466 & 0.918 & 0.347 \\
arvalis~1   & 827 & 0.403 & 0.890 & 0.369 \\
usask~1     & 153 & 0.408 & 0.874 & 0.384 \\
arvalis~2   & 162 & 0.365 & 0.813 & 0.431 \\
\bottomrule
\end{tabular}%
}
\end{table}

% ----------------------------------------------------------
\section{Augmentation Experiments}

\subsection{Setup}

All models use YOLOv8s initialized from COCO pretrained weights.
%
% [FILL] Insert exact dataset split counts, e.g.:
% MinneApple: 536 train / 67 val / 67 test images,
%             ~23,000 annotated instances.
% GWHD: 2,700 train / 735 val / 748 test images.
%
Training: MinneApple 200 epochs / patience 10; GWHD 100 epochs /
patience 8; batch 16; input $1024{\times}1024$; lr $10^{-2}$.
All Ultralytics built-in augmentations are disabled to isolate
copy-paste effects. Augmentation ratio 0.3; scale jitter
$[0.9,1.1]$ (MA) / $[0.8,1.2]$ (GWHD). Results are 3-seed means.

\subsection{Results and Feature-Space Analysis}

Table~\ref{tab:ap} shows that on MinneApple, score-guided copy-paste
outperforms random copy-paste ($\Delta\mathrm{AP}{=}{+}0.013$) across
all sub-metrics, but both conditions fall below baseline
($-0.024$ for Score CP). The AP$_{75}$ drop ($0.452{\to}0.440$)
indicates degraded localization precision, consistent with hard objects
pasted without blending or masks, causing the model to learn edge
artifacts. The AP$_S$ drop ($0.300{\to}0.286$) reflects that
score-selected hard objects are predominantly sub-resolution---they
cannot be made more learnable regardless of paste frequency.
On GWHD, score-guided and random are indistinguishable ($-0.002$),
and both degrade below baseline ($-0.023$). Score-guided sampling
concentrates objects from domain-minority sources (\textit{arvalis}),
amplifying existing domain imbalance.

\begin{table}[t]
\centering
\caption{Three-seed mean COCO AP. \textbf{Bold} = best copy-paste
  condition per dataset.}
\label{tab:ap}
\resizebox{\columnwidth}{!}{%
\begin{tabular}{@{}llrrrrr@{}}
\toprule
\textbf{Dataset} & \textbf{Condition}
  & \textbf{AP}
  & \textbf{AP\textsubscript{50}}
  & \textbf{AP\textsubscript{75}}
  & \textbf{AP\textsubscript{S}}
  & \textbf{AP\textsubscript{M}} \\
\midrule
\multirow{3}{*}{MinneApple}
  & Baseline          & 0.439 & 0.771 & 0.452 & 0.300 & 0.592 \\
  & Random CP         & 0.415 & 0.741 & 0.427 & 0.279 & 0.563 \\
  & \textbf{Score CP} & \textbf{0.428} & \textbf{0.759}
                      & \textbf{0.440} & \textbf{0.286} & \textbf{0.581} \\
\midrule
\multirow{3}{*}{GWHD}
  & Baseline          & 0.507 & 0.922 & 0.491 & 0.136 & 0.501 \\
  & Random CP         & 0.486 & 0.905 & 0.456 & 0.118 & 0.480 \\
  & \textbf{Score CP} & \textbf{0.484} & \textbf{0.905}
                      & \textbf{0.455} & \textbf{0.098} & \textbf{0.477} \\
\bottomrule
\end{tabular}%
}
\end{table}

% ----------------------------------------------------------------
% FIGURE 2 RECOMMENDATION:
% Replace this placeholder with a 2-panel UMAP plot:
%   Left: MinneApple backbone features (P4, global avg pool).
%   Right: GWHD backbone features.
% Four point groups per panel, different colors/markers:
%   - Original train images
%   - Test images
%   - Random CP augmented images
%   - Score CP augmented images
% Mark group centroids with a large X.
% This figure carries the distributional shift argument.
% Export as umap_combined.pdf (width = \columnwidth).
% ----------------------------------------------------------------
\begin{figure}[t]
\centering
%\includegraphics[width=\columnwidth]{umap_combined.pdf}
\framebox{\parbox{0.88\columnwidth}{\centering\vspace{1.4cm}
\small FIGURE 2 PLACEHOLDER\\
UMAP of YOLOv8s P4 backbone features.\\
Left: MinneApple \quad Right: GWHD\\
Groups: original train, test, random CP, score CP.
\vspace{1.4cm}}}
\caption{Image-level backbone feature projections (YOLOv8s P4, UMAP).
  Augmented images are displaced from the test distribution relative to
  original training images on both datasets.
  % [FILL] Add centroid distance values, e.g.:
  % MinneApple: Score CP centroid distance to test = X.XX vs.
  %             original train = X.XX.
  On GWHD, Score CP shows greater displacement than Random CP,
  consistent with domain-minority object concentration.}
\label{fig:umap}
\end{figure}

% ----------------------------------------------------------
\section{Discussion}

\textbf{Score validity vs.\ augmentation utility are separable claims.}
The validation results (Section~III) establish that MinImageScorer
correctly identifies hard objects and images: domain difficulty
rankings are recovered without test labels, object size difficulty
tracks ground-truth AP, and per-image correlations with miss rate
reach $r{=}0.82$ on GWHD. This is the primary contribution.

\textbf{Why copy-paste fails here.}
Kisantal~\textit{et~al.}~\cite{kisantal2019augmentation} and
Ghiasi~\textit{et~al.}~\cite{ghiasi2021simple} demonstrate copy-paste
gains on MS COCO, where hard objects are above the resolution floor
and within-distribution. Neither condition holds on MinneApple
(sub-resolution hard objects) or GWHD (cross-domain images). The score
selects the right objects; the method cannot recover from a structural
limitation that precedes selection.

\textbf{Practical implication.}
Run MinImageScorer before applying copy-paste: if high-score objects
are predominantly sub-resolution or cross-domain, copy-paste will
degrade performance regardless of selection quality. This constitutes
a dataset-level diagnostic criterion absent from prior work.

% ----------------------------------------------------------
\section{Conclusion}

MinImageScorer reliably quantifies object and image detection
difficulty from model failure signals, validated by strong correlation
with miss rate and FP rate, and by correctly recovering domain- and
size-level difficulty rankings without supervision.
Score-guided copy-paste outperforms random selection on MinneApple
($+0.013$~AP) but both strategies degrade below baseline, explained
by distributional shift in feature space. On GWHD, domain shift
renders copy-paste ineffective under any selection strategy.
Future work includes online augmentation to reduce distributional
rigidity and domain-constrained pasting for multi-source benchmarks.

% ----------------------------------------------------------
\begin{thebibliography}{00}

\bibitem{hani2020minneapple}
N.~Hani, P.~Roy, and B.~Bhanu,
``MinneApple: A Benchmark Dataset for Apple Detection and Segmentation,''
\textit{IEEE Robotics and Automation Letters},
vol.~5, no.~2, pp.~852--858, 2020.

\bibitem{david2021global}
E.~David \textit{et~al.},
``Global Wheat Head Detection (GWHD) Dataset,''
\textit{Plant Phenomics}, vol.~2021, 2021.

\bibitem{kisantal2019augmentation}
M.~Kisantal \textit{et~al.},
``Augmentation for Small Object Detection,''
in \textit{Proc.\ 9th Int.\ Conf.\ Advances in Computing,
Communication and Informatics (ICACCI)}, 2019.

\bibitem{ghiasi2021simple}
G.~Ghiasi \textit{et~al.},
``Simple Copy-Paste is a Strong Data Augmentation Method
for Instance Segmentation,''
in \textit{Proc.\ IEEE/CVF CVPR}, pp.~2918--2928, 2021.

\bibitem{yolov8}
G.~Jocher \textit{et~al.},
``Ultralytics YOLOv8,''
GitHub, 2023. [Online].
Available: \texttt{https://github.com/ultralytics/ultralytics}

\end{thebibliography}

\end{document}
